% V. Методика за експериментална проверка.
% Разделът излага ПЛАНА на експериментите. Не съдържа нито едно число:
% всички стойности се попълват от измерванията в Раздел VI.
\section{Методика за експериментална проверка}
\label{sec:method}

Разделът описва как се получават резултатите от Раздел~\ref{sec:results}. Той се
пише преди измерванията, а не след тях, за да може погрешно очакване да остане
видимо и да бъде отчетено като погрешно.

\subsection{Среда за провеждане}
\label{subsec:environment}

Всички опити се провеждат в изолирана среда, изградена изцяло от автора. Клиентът и
сървърът работят върху локален интерфейс или в контейнери на една машина. Никакъв
трафик не напуска тази среда и никаква външна услуга не участва в опитите. Това
ограничение е задължително и не подлежи на облекчаване. Точната конфигурация на
машината и на средата се описва в \TODO{конфигурация на средата: процесор, памет,
ядро, версии на инструментите}.

\subsection{Набори от данни}
\label{subsec:datasets}

Анализаторът на ръкостискането се проверява върху записани файлове в PCAP формат,
получени от собствената среда. Всеки запис се придружава от очаквания резултат,
записан ръчно след разчитане на съответната спецификация. Проверката на веригата се
провежда върху собствена йерархия от сертификати, изградена за целта, с нарочно
въведени дефекти по един на сертификат: изтекъл срок, липсващо основно ограничение,
несъответствие на името, отменен сертификат. Пълният списък на случаите се дава в
\TODO{таблица със случаите за проверка на веригата и очаквания резултат за всеки}.

\subsection{План на измерванията при лавинен трафик}
\label{subsec:load_plan}

Измерва се цената на всеки механизъм за допускане поотделно и на съчетанията им.
За всяка конфигурация се снема закъснението при установяване на връзка, броят на
успешно установените връзки за единица време и потребената памет за състояние. Всяка
конфигурация се повтаря многократно, а докладваната стойност е медиана заедно с
разсейването, а не единично измерване.

Измерва се и базова конфигурация, при която нито един механизъм не е включен. Без
нея числата за останалите конфигурации нямат отправна точка. Броят на повторенията,
продължителността на един опит и загряващият период се определят предварително и се
записват в \TODO{параметри на опита: брой повторения, продължителност, загряване}.

\subsection{Изисквания за валидност}
\label{subsec:validity}

Измерването се приема за валидно само ако са изпълнени следните условия. Генераторът
на натоварване не трябва да бъде тясното място, което се проверява с отделен опит
срещу нищо неправещ приемник. Машината не трябва да е ограничена термично или от
режима на управление на честотата. Разсейването между повторенията трябва да остане
под предварително определен праг, иначе резултатът се отхвърля, а не се осреднява.
Праговете се задават в \TODO{числени прагове за приемане на измерването}.

\subsection{Очаквани резултати}
\label{subsec:hypotheses}

Очакванията се записват тук преди измерването. Очаква се механизмът без съхранение
на състояние да намали потребената памет при лавина за сметка на ограничение върху
договаряните опции. Очаква се ограничаването по източник да има пренебрежимо
закъснение и слаб ефект при разпределен източник. Очаква се допускането срещу
изчислителна задача да внесе най-голямо закъснение при първо свързване. Дали тези
очаквания се потвърждават, се отчита в Раздел~\ref{sec:results}, включително когато
не се потвърждават.
